% T3a: Morphological coherence + subword ablation by frequency band
% (paper_experiments.py, analysis "morph"; related vs frequency-matched control)
\begin{table}[t]
  \centering
  \begin{tabular}{llrrrrr}
    \toprule
    \textbf{Representation} & \textbf{Band} & \textbf{$n$} &
    \textbf{related} & \textbf{control} & \textbf{margin} & \textbf{acc} \\
    \midrule
    \multirow{2}{*}{FastText (char $n$-gram)}
      & freq $\geq 100$ & 27 & 0.601 & 0.284 & $+0.317$ & 0.96 \\
      & rare $< 100$    & 22 & 0.640 & 0.353 & $+0.287$ & 0.95 \\
    \addlinespace
    \multirow{2}{*}{word2vec (no subword)}
      & freq $\geq 100$ & 27 & 0.490 & 0.365 & $+0.125$ & 0.85 \\
      & rare $< 100$    & 22 & 0.498 & 0.422 & $+0.076$ & 0.77 \\
    \addlinespace
    \multirow{2}{*}{byte-LM}
      & freq $\geq 100$ & 27 & 0.988 & 0.967 & $+0.021$ & 0.93 \\
      & rare $< 100$    & 22 & 0.992 & 0.968 & $+0.025$ & 1.00 \\
    \addlinespace
    \multirow{2}{*}{char-Transformer}
      & freq $\geq 100$ & 27 & 0.916 & 0.790 & $+0.126$ & 0.96 \\
      & rare $< 100$    & 22 & 0.936 & 0.774 & $+0.162$ & 1.00 \\
    \bottomrule
  \end{tabular}
  \caption{Morphological coherence by frequency band: mean cosine between
  root-sharing forms (\emph{related}) versus a frequency-matched form on a
  different root (\emph{control}), over 8 curated triconsonantal root families
  (49 related pairs). FastText shows the largest margin; the no-subword
  \textsc{word2vec} baseline is weaker and \emph{degrades for rare forms}
  ($+0.125\!\to\!+0.076$), whereas FastText stays robust. Neural models'
  raw cosines are compressed by anisotropy but still rank related above control.}
  \label{tab:morph}
\end{table}
